\documentclass[12pt]{article}
\usepackage[utf8]{inputenc}
\usepackage[T1]{fontenc}
\usepackage[margin=1in]{geometry}
\usepackage{setspace}
\usepackage{xcolor}
\usepackage[colorlinks=true, urlcolor=blue!60!black, linkcolor=blue!60!black]{hyperref}
\setlength{\parindent}{2em}
\setlength{\parskip}{1.25em}
\renewcommand{\baselinestretch}{1.0}
\begin{document}
\begin{center}
{\LARGE \textit{Wisconsin v. Mitchell}, 508 U.S. 476 (1993)\par}
\vspace{0.8em}
{\large Case Brief --- Four-Page Edition\par}
\end{center}
\vspace{1em}

\section{Memory Jogger}

Unlike laws that criminalize bias-motivated speech, a hate-crime penalty enhancement that increases the punishment for unprotected criminal conduct when the defendant selects the victim because of race or another protected characteristic does not violate the First Amendment (\href{https://www.courtlistener.com/opinion/112879/wisconsin-v-mitchell/}{CourtListener}).

\section{Detailed Case Facts}

On October 7, 1989, Todd Mitchell, a young black man, was with a group of young black men and boys at an apartment complex in Kenosha, Wisconsin. Some discussed a scene from \textit{Mississippi Burning} in which a white man beat a young black boy who was praying. Outside, Mitchell asked, ``Do you all feel hyped up to move on some white people?'' When a young white boy walked by on the other side of the street, Mitchell pointed him out, counted to three, and sent the group after him. They beat the boy severely and stole his tennis shoes; he was in a coma for four days.

Mitchell was charged with aggravated battery, which ordinarily carried a two-year maximum. Prosecutors also invoked Wis. Stat. \S~939.645, which raised the maximum for an underlying crime by up to five years when the defendant ``[i]ntentionally selects the person against whom the crime \ldots\ is committed \ldots\ because of the race, religion, color, disability, sexual orientation, national origin or ancestry of that person.''

\section{Procedural History}

A Kenosha County jury convicted Mitchell and found by special verdict that he had selected the victim because of race, raising his maximum from two to seven years; he was sentenced to four years. The Wisconsin Court of Appeals affirmed. The Wisconsin Supreme Court reversed, holding under \textit{R.A.V. v. City of St. Paul} that the statute punished ``offensive thought'' and was overbroad because using a defendant's prior speech to prove motive would chill expression. The U.S. Supreme Court granted certiorari to resolve a split among state high courts, heard argument on April 21, 1993, and decided the case on June 11, 1993.

\section{Judicial Votes}

Reversed and remanded, \textbf{9--0}.
\begin{itemize}
\item \textbf{Majority:} Chief Justice Rehnquist, joined by Justices White, Blackmun, Stevens, O'Connor, Scalia, Kennedy, Souter, and Thomas.
\item \textbf{Concurrences:} None.
\item \textbf{Dissents:} None.
\end{itemize}

\section{Holding}

The First Amendment does not prohibit a State from enhancing the penalty for unprotected criminal conduct when the defendant intentionally selects the victim because of race or other protected status, nor does it bar using a defendant's prior speech to prove motive or intent, subject to ordinary rules of evidence. The Court was not bound by the Wisconsin Supreme Court's characterization of the statute's practical effect. Its observations about the societal harms of bias crimes are supporting reasoning, and it declined to reach Mitchell's equal protection and vagueness claims.

\section{Analysis of Opinions}

\textbf{Not bound by the state court.} Federal courts accept a state court's construction of a statute's words, but the Wisconsin court had characterized the law's ``practical effect,'' a federal question.

\textbf{Conduct and motive.} The same battery drew a harsher maximum only because of motive, but physical assault is not expressive conduct. Judges have long weighed motive, such as murder for pecuniary gain. \textit{Dawson v. Delaware} forbids punishing ``abstract beliefs'' unrelated to the crime, but \textit{Barclay v. Florida} permits considering racial animus tied to the offense, even in capital sentencing.

\textbf{Antidiscrimination laws and \textit{R.A.V.}} Motive plays the same role as under Title VII, which prohibits conduct ``because of'' race and has survived First Amendment challenge. The St. Paul ordinance in \textit{R.A.V.} was directed at expression; Wisconsin's statute targets unprotected conduct. Bias crimes are also reasonably thought to inflict distinct emotional harm and provoke retaliation and community unrest, a justification beyond disagreement with bigoted beliefs.

\textbf{Overbreadth and evidence.} The chilling theory was too speculative: citizens are unlikely to suppress bigoted views for fear they might someday be offered at a felony trial. Under \textit{Haupt} and \textit{Price Waterhouse}, prior statements may prove motive, subject to relevance and reliability. No Justice wrote separately.

\section{Examples: Future Applications}

\begin{enumerate}
\item \textbf{Same-Side (Constitutional): Church Arson.} Enhancing an arson sentence for burning a house of worship because of the congregation's religion, with the defendant's social-media posts admitted to prove motive, is valid.
\item \textbf{Same-Side (Constitutional): Vandalism.} Enhancing vandalism penalties when property is targeted because of the owner's sexual orientation or disability mirrors \S~939.645.
\item \textbf{Opposite-Side (Unconstitutional): Speech Ordinance.} Jailing anyone who displays symbols or shouts epithets that arouse anger ``on the basis of race, religion, or gender,'' with no underlying crime, falls under \textit{R.A.V.}
\item \textbf{Opposite-Side (Unconstitutional): Unrelated Beliefs.} Using a burglar's extremist membership and racist views to increase a sentence for a non-bias burglary violates \textit{Dawson}.
\item \textbf{Fence-Sitter: Disorderly Conduct at a Rally.} Enhancing a disorderly-conduct penalty for a protester who screams slurs and shoves a sign depends on whether the conviction rests on speech or on physical conduct.
\end{enumerate}

\section{Critique}

\textbf{Civil-libertarian critiques.} Susan Gellman, who co-authored Mitchell's brief, Steven Gey, and others argue that because the battery is already punished, the only trigger for five more years is the defendant's ideological motive. They see tension with \textit{R.A.V.}, since a legislature can evade strict scrutiny by attaching a bias penalty to an existing crime, and James Jacobs and Kimberly Potter question whether enhancements deter bias violence.

\textbf{Progressive and civil-rights views.} Frederick Lawrence and groups including the ACLU and the Anti-Defamation League praise the decision, stressing that bias violence terrorizes whole communities and undermines equal citizenship. Grounded in sentencing tradition and Title VII, the unanimous ruling gave legislatures a workable blueprint without eroding protection for offensive speech.

\section{Key Quotations}

\begin{quote}
``The First Amendment, moreover, does not prohibit the evidentiary use of speech to establish the elements of a crime or to prove motive or intent.''
\end{quote}
\noindent --- Chief Justice Rehnquist, Opinion of the Court, 508 U.S. at 489.

\begin{quote}
``Thus, a physical assault is not by any stretch of the imagination expressive conduct protected by the First Amendment.''
\end{quote}
\noindent --- Chief Justice Rehnquist, Opinion of the Court, 508 U.S. at 484.

\begin{quote}
``Moreover, the Wisconsin statute singles out for enhancement bias-inspired conduct because this conduct is thought to inflict greater individual and societal harm.''
\end{quote}
\noindent --- Chief Justice Rehnquist, Opinion of the Court, 508 U.S. at 487--488.

\end{document}
